 \documentclass[12pt,reqno]{amsart}
%%%%%%%%%%%%%%%%%%%%%%%%%%%%%%%%%%%%%%%%%%%%%%%%%%%
%								Packages									         %
%%%%%%%%%%%%%%%%%%%%%%%%%%%%%%%%%%%%%%%%%%%%%%%%%%%
\usepackage[T1]{fontenc}
\usepackage{amsmath, amssymb, amsthm, amscd, amsfonts}
\usepackage{mathtools}
\usepackage{enumerate,multicol}
\usepackage[shortlabels]{enumitem}

\usepackage[dvipsnames]{xcolor}
\definecolor{DartmouthGreen}{HTML}{00693E}
\definecolor{DartmouthDarkGreen}{HTML}{004B23}
\definecolor{DartmouthLightGreen}{HTML}{4BAF7A}

\usepackage[
  colorlinks=true,
  hyperindex,
  linkcolor=DartmouthDarkGreen,
  citecolor=DartmouthGreen,
  urlcolor=DartmouthGreen,
  pagebackref=true
]{hyperref}

\usepackage{parskip}
\usepackage[letterpaper, margin=.5in, top=1in, bottom=1.4in, headheight=42pt, headsep=14pt, footskip=50pt]{geometry}
\linespread{1.1}

\usepackage{algorithm,algpseudocode,verbatim}
\usepackage{float,caption,comment}

\usepackage{tikz,tikz-cd}
\usetikzlibrary{graphs, graphs.standard,positioning}



\usepackage{calrsfs}
\DeclareMathAlphabet{\pazocal}{OMS}{zplm}{m}{n}
\usepackage{newcent,fouriernc}

%%%%%%%%%%%%%%%%%%%%%%%%%%%%%%%%%%%%%%%%%%%%%%%%%%%
%								Copyright 								         %
%%%%%%%%%%%%%%%%%%%%%%%%%%%%%%%%%%%%%%%%%%%%%%%%%%%
\usepackage{ccicons}
\usepackage{fancyhdr}

\pagestyle{fancy}
\fancyhf{}
\fancyfoot[R]{\footnotesize \thepage}
\renewcommand{\headrulewidth}{0pt}
\renewcommand{\footrulewidth}{0pt}

\newcommand{\originalurl}{}
\newcommand{\originalurlI}{}

\fancypagestyle{firstpage}{%
  \fancyhf{}
\fancyhead[L]{\footnotesize Dartmouth College \\ Math 121: Commutative Algebra}
\fancyhead[R]{\footnotesize \href{https://www.juliettebruce.xyz/teaching/121S26}{Spring 2026}}
\fancyfoot[L]{\footnotesize \ccbyncsa\ \ \ccCopy \ Juliette Bruce 2026.\\
Licensed under \href{https://creativecommons.org/licenses/by-nc-sa/4.0/}{Creative Commons BY-NC-SA 4.0}.\\
Adapted from \ccCopy \href{https://sites.lsa.umich.edu/kesmith/}{Karen Smith} (2019); see \href{\originalurl}{original}.}
\fancyfoot[R]{\footnotesize \thepage}
\renewcommand{\headrulewidth}{0.4pt}
\renewcommand{\footrulewidth}{0.4pt}
}

\fancypagestyle{firstpage2}{%
  \fancyhf{}
\fancyhead[L]{\footnotesize Dartmouth College \\ Math 121: Commutative Algebra}
\fancyhead[R]{\footnotesize \href{https://www.juliettebruce.xyz/teaching/121S26}{Spring 2026}}
\fancyfoot[L]{\footnotesize \ccbyncsa\ \ \ccCopy \ Juliette Bruce 2026.\\
Licensed under \href{https://creativecommons.org/licenses/by-nc-sa/4.0/}{Creative Commons BY-NC-SA 4.0}.\\
Adapted from \ccCopy \href{https://sites.lsa.umich.edu/kesmith/}{Karen Smith} (2019); see \href{\originalurl}{original} and \href{\originalurlI}{original}.}
\fancyfoot[R]{\footnotesize \thepage}
\renewcommand{\headrulewidth}{0.4pt}
\renewcommand{\footrulewidth}{0.4pt}
}


\fancypagestyle{firstpageIP}{%
  \fancyhf{}
\fancyhead[L]{\footnotesize Dartmouth College \\ Math 121: Commutative Algebra}
\fancyhead[R]{\footnotesize \href{https://www.juliettebruce.xyz/teaching/121S26}{Spring 2026}}
\fancyfoot[L]{\footnotesize \ccbyncsa\ \ \ccCopy \ Juliette Bruce 2026.\\
Licensed under \href{https://creativecommons.org/licenses/by-nc-sa/4.0/}{Creative Commons BY-NC-SA 4.0}.\\
Inspired by \ccCopy\ \href{https://sites.lsa.umich.edu/kesmith/}{Karen Smith} (2019); \href{https://sites.lsa.umich.edu/kesmith/teaching/math-614-commutative-algebra/}{see}.}
\fancyfoot[R]{\footnotesize \thepage}
\renewcommand{\headrulewidth}{0.4pt}
\renewcommand{\footrulewidth}{0.4pt}
}
%%%%%%%%%%%%%%%%%%%%%%%%%%%%%%%%%%%%%%%%%%%%%%%%%%%
%								Theorems 								         %
%%%%%%%%%%%%%%%%%%%%%%%%%%%%%%%%%%%%%%%%%%%%%%%%%%%

\newtheorem{maintheorem}{Theorem}
\newtheorem{theorem}{Theorem}
\newtheorem{lemma}[theorem]{Lemma}

\newtheorem{corollary}[theorem]{Corollary}
\newtheorem{proposition}[theorem]{Proposition}
\newtheorem{conjecture}[theorem]{Conjecture}


\theoremstyle{definition}
\newtheorem{maindefinition}[maintheorem]{Definition}
\newtheorem{definition}[theorem]{Definition}
\newtheorem{setup}[theorem]{Set-Up}
\newtheorem{notation}[theorem]{Notation}
\newtheorem{convention}[theorem]{Convention}
\newtheorem{construction}[theorem]{Construction}





\setcounter{MaxMatrixCols}{18}

%%%%%%%%%%%%%%%%%%%%%%%%%%%%%%%%%%%%%%%%%%%%%%%%%%%
%								Letters									         %
%%%%%%%%%%%%%%%%%%%%%%%%%%%%%%%%%%%%%%%%%%%%%%%%%%%


\newcommand{\CC}{\mathbb{C}}
\newcommand{\FF}{\mathbb{F}}
\newcommand{\EE}{\mathbb{E}}

\newcommand{\GG}{\mathbb{G}}
\newcommand{\HH}{\mathbb{H}}
\newcommand{\II}{\mathbb{I}}
\newcommand{\KK}{\mathbb{K}}

\newcommand{\MM}{\mathbb{M}}
\newcommand{\NN}{\mathbb{N}}
\newcommand{\PP}{\mathbb{P}}
\newcommand{\QQ}{\mathbb{Q}}
\newcommand{\RR}{\mathbb{R}}
\renewcommand{\SS}{\mathbb{S}}
\newcommand{\TT}{\mathbb{T}}
\newcommand{\VV}{\mathbb{V}}
\newcommand{\WW}{\mathbb{W}}

\newcommand{\ZZ}{\mathbb{Z}}


\newcommand{\cA}{\mathcal{A}}
\newcommand{\cB}{\mathcal{B}}
\newcommand{\cC}{\mathcal{C}}
\newcommand{\cE}{\mathcal{E}}

\newcommand{\cF}{\mathcal{F}}
\newcommand{\cG}{\mathcal{G}}

\newcommand{\cH}{\mathcal{H}}
\newcommand{\cI}{\mathcal{I}}
\newcommand{\cL}{\mathcal{L}}
\newcommand{\cM}{\mathcal{M}}
\newcommand{\cO}{\mathcal{O}}
\newcommand{\cP}{\mathcal{P}}
\newcommand{\cQ}{\mathcal{Q}}
\newcommand{\cS}{\mathcal{S}}
\newcommand{\cT}{\mathcal{T}}
\newcommand{\cR}{\mathcal{R}}
\newcommand{\cU}{\mathcal{U}}
\newcommand{\cV}{\mathcal{V}}

\newcommand{\pB}{\pazocal{B}}
\newcommand{\pC}{\pazocal{C}}
\newcommand{\pO}{\pazocal{O}}
\newcommand{\pG}{\pazocal{G}}

\newcommand{\sM}{\mathsf{M}}
\newcommand{\sN}{\mathsf{N}}

\newcommand{\sQ}{\mathsf{Q}}
\newcommand{\sP}{\mathsf{P}}

\newcommand{\fm}{\mathfrak{m}}
\newcommand{\fp}{\mathfrak{p}}
\newcommand{\fq}{\mathfrak{q}}

%%%%%%%%%%%%%%%%%%%%%%%%%%%%%%%%%%%%%%%%%%%%%%%%%%%
%								Operators									         %
%%%%%%%%%%%%%%%%%%%%%%%%%%%%%%%%%%%%%%%%%%%%%%%%%%%
\DeclareMathOperator{\Id}{Id}
\DeclareMathOperator{\ev}{ev}
\DeclareMathOperator{\ord}{ord}
%
\DeclareMathOperator{\Hom}{Hom}
\DeclareMathOperator{\End}{End}
\DeclareMathOperator{\coker}{coker}
\DeclareMathOperator{\Ann}{Ann}
\DeclareMathOperator{\img}{img}
%
\DeclareMathOperator{\Frac}{Frac}
\DeclareMathOperator{\Nil}{Nil}
\DeclareMathOperator{\red}{red}
%
\DeclareMathOperator{\Spec}{Spec}
\DeclareMathOperator{\mSpec}{mSpec}
%
\DeclareMathOperator{\SL}{SL}
%
\DeclareMathOperator{\init}{in}
\DeclareMathOperator{\lex}{lex}
\DeclareMathOperator{\grlex}{grlex}
\DeclareMathOperator{\grevlex}{grevlex}
\DeclareMathOperator{\revlex}{revlex}
\DeclareMathOperator{\LT}{LT}
\DeclareMathOperator{\LM}{LM}
\DeclareMathOperator{\LC}{LC}


%%%%%%%%%%%%%%%%%%%%%%%%%%%%%%%%%%%%%%%%%%%%%%%%%%%
%								Categories									         %
%%%%%%%%%%%%%%%%%%%%%%%%%%%%%%%%%%%%%%%%%%%%%%%%%%%

\DeclareMathOperator{\comRing}{\textbf{comRing}}
\DeclareMathOperator{\Top}{\textbf{Top}}
\DeclareMathOperator{\Set}{\textbf{Set}}
\DeclareMathOperator{\alg}{\textbf{alg}}
\DeclareMathOperator{\Mod}{\textbf{Mod}}


%%%%%%%%%%%%%%%%%%%%%%%%%%%%%%%%%%%%%%%%%%%%%%%%%%%
%								MISC									         %
%%%%%%%%%%%%%%%%%%%%%%%%%%%%%%%%%%%%%%%%%%%%%%%%%%%
\makeatletter
\newcommand*{\tp}{%
  {\mathpalette\@transpose{}}%
}
\newcommand*{\@transpose}[2]{%
  % #1: math style
  % #2: unused
  \raisebox{\depth}{$\m@th#1\intercal$}%
}
\makeatother

\newcommand{\juliette}[1]{{\color{purple} \sf  Juliette: [#1]}}
%%%%%%%%%%%%%%%%%%%%%%%%%%%%%%%%%%%%%%%%%%%%%%%%%%%
%%%%%%%%%%%%%%%%%%%%%%%%%%%%%%%%%%%%%%%%%%%%%%%%%%%
%%%%%%%%%%%%%%%%%%%%%%%%%%%%%%%%%%%%%%%%%%%%%%%%%%%
%%%%%%%%%%%%%%%%%%%%%%%%%%%%%%%%%%%%%%%%%%%%%%%%%%%
\title{Worksheet 1.1: Basics of Rings}

%%%%%%%%%%%%%%%%%%%%%%%%%%%%%%%%%%%%%%%%%%%%%%%%%%%
%%%%%%%%%%%%%%%%%%%%%%%%%%%%%%%%%%%%%%%%%%%%%%%%%%%


%%%%%%%%%%%%%%%%%%%%%%%%%%%%%%%%%%%%%%%%%%%%%%%%%%%
%%%%%%%%%%%%%%%%%%%%%%%%%%%%%%%%%%%%%%%%%%%%%%%%%%%

\begin{document}

%%%%%%%%%%%%%%%%%%%%%%%%%%%%%%%%%%%%%%%%%%%%%%%%%%%
%%%%%%%%%%%%%%%%%%%%%%%%%%%%%%%%%%%%%%%%%%%%%%%%%%%

\maketitle
\thispagestyle{firstpageIP}

A \emph{ring} $(R,+,\cdot)$ is an abelian group $(R,+)$ together with a multiplication $(\cdot)$ and a multiplicative identity element $1_{R} \in R$ satisfying the following axioms:
\begin{enumerate}[(i)]
	\item (Multiplicative Identity): $1_{R} \cdot a = a \cdot 1_{R} = a$ for all $a\in R$,
	\item (Associativity): $a\cdot (b\cdot c) = (a\cdot b) \cdot c$ for all $a,b,c\in R$, and
	\item (Distributivity): $a\cdot(b+c) = (a\cdot b) + (a\cdot c)$ and $(b+c) \cdot a = (b\cdot a) + (c\cdot a)$ for all $a,b,c \in R$.
\end{enumerate}
As a warning some people will not include condition (i) in the definition of a ring, but this is rare to unheard of in modern commutative algebra. (See the historical remark at the end of this worksheet.) As is often the case once we are familiar with these operations we will normally just drop some of the notational baggage and just write $ab$ for $a\cdot b$ and $1$ for $1_{R}$. 

A ring $R$ is \emph{commutative} if and only if $ab = ba$ for all $a,b \in R$. As the name of the course suggests we will primarily be interested in commutative rings in this course. To the point we shall adopt the following convention: Throughout this course, "ring" means  \textit{commutative} ring with unity. Commutative algebra broadly can be thought of as the study of rings and their associated objects. At a high level we will study rings from four perspectives: 1) element-wise, 2) homomorphisms, 3) ideals, and 4) modules. You hopefully have some familiarity with 1, 2, and 3, but today we will review and remember what it is like working with these objects. 

An element $u \in R$ is a \emph{unit} if there exists $v \in R$ such that $uv = vu = 1$, in which case we will often denote $v$ by $u^{-1}$. A nonzero element $a \in R$ is a \emph{zero divisor} if there exists a nonzero element $b\in R$ such that $ab=0$. A non-zero element, which is not a zero divisor is a \emph{non-zero divisor}. 
 \hrulefill 
	
\begin{enumerate}[leftmargin=*]
	 \item Let $R$ be a ring and consider $R[[t]]$ the ring of formal power series with coefficients in $R$. The goal of this exercise is to prove that:
	 \[
	 \left(R[[t]]\right)^{\times} = \left\{ u + tf(t) \;\; \big| \;\; u \in R^{\times}, \;\; f(t) \in R[[t]] \right\}.
	 \]
	 \begin{enumerate}[label=(\alph*),ref=\theenumi\alph*]
	 	\item\label{ex-pt:power-series-ct-term} Elements of $R[[t]]$ can be written as $f=\sum_{k=0}^{\infty} f_k t^k$ for $f_{k} \in R$. We call $f_0$ the constant term of $f$. If $f=\sum_{k=0}^{\infty} f_k t^k$ and $g=\sum_{k=0}^{\infty} g_k t^k$ then find a formula for the constant term of $fg$. 
		\item Using the formula from \ref{ex-pt:power-series-ct-term} prove that if $f$ is a unit then $f_0 \in R^{\times}$. 
		\item\label{ex-pt:power-series-mult} Find a formula for the coefficient of $t^k$ appearing in $fg$. 
		\item Assume that $f_0 \in R^{\times}$. Prove that $f$ is a unit by explicitly constructing an element $g \in R[[t]]$ such that $fg = 1$. (Hint: The condition that $fg = 1$ places constraints on each coefficient of $fg$. Use the constraints together with part \ref{ex-pt:power-series-mult} to recursively solve for the coefficients of $g$.)	\end{enumerate}	
\end{enumerate}	

\hrulefill

A non-empty subset $I \subset R$ is an \emph{ideal} if it is closed by outside multiplication by elements of $R$ and is closed under addition. Precisely, $I\subset R$ is a non-empty subset satisfying:
\begin{enumerate}[(i)]
	\item if $a \in I$ then $ra \in I$ for all $r \in R$ and
	\item if $a,b \in I$ then $a+b \in I$. 
\end{enumerate}
Note this immediately implies that an ideal always contains $0$ and is an abelian group with respect to addition. Given elements $a_1,\ldots,a_t \in R$ we write $\langle a_1,\ldots, a_t\rangle \subset R$ for the the smallest ideal containing $a_1,\ldots,a_t$. Concretely, we may describe this as the set of all $R$-linear combinations of $a_1,\ldots, a_t$:
\[
\langle a_1,\ldots,a_t \rangle = \left\{ r_1a_1+\cdots +r_ta_t \;\; \big|\;\; r_1,\ldots,r_t\in R\right\}.
\]
An ideal $I \subset R$ is \emph{principal} if there exists an element $a\in R$ such that $I = \langle a \rangle$. 

 \hrulefill 

\begin{enumerate}[leftmargin=*]
\item 
	\begin{enumerate}[label=(\alph*),ref=\theenumi\alph*]
		\item Prove that an element $u \in R$ is a unit if and only if the principal ideal $\langle u \rangle = R$.
		\item Prove that if $ab$ and $b$ are units, then $a$ is a unit.
		\item Prove that if $u$ is a unit and $a \in R$, then $a$ and $ua$ generate the same principal ideal.
		\item In the ring $\ZZ/\langle 12 \rangle$, determine which elements are units.
	\end{enumerate}
\end{enumerate}

 \hrulefill 

Let $R$ and $S$ be rings. A ring homomorphism from $R$ to $S$ is a function $\phi:R \to S$ such that:
\begin{enumerate}[(i)]
	\item (Multiplicative): $\phi(ab) = \phi(a)\phi(b)$ for all $a,b\in R$,
	\item (Additive): $\phi(a+b) = \phi(a)+\phi(b)$ for all $a,b\in R$, and
	\item (Preserves Unit): $\phi(1_R) = 1_S$.
\end{enumerate}
Notice that condition (ii) implies that $\phi(0_R) = 0_S$, so $\phi$ is a homomorphism of the underlying abelian groups. It will turn out that condition (iii) is very useful as the following exercise demonstrates. There are two important sets associated to any ring homomorphism $\phi:R \to S$ the kernel and image of $\phi$:
\[
\ker(\phi) \coloneqq \left\{ a \in R \;\; \big| \; \; \phi(a) = 0 \right\}\subset R \quad \quad \text{and} \quad \quad \img(\phi) \coloneqq \left\{ \phi(a) \in S \;\; \big| \; \; a\in R\right\} \subset S.
\]
More succinctly, $\ker(\phi) = \phi^{-1}(\{0\})$ and $\img(\phi) = \phi(R)$. One checks that $\ker(\phi)$ is an ideal of $R$, and that $\img(\phi)$ is a subring of $S$; however, it need not be an ideal of $S$. 

 \hrulefill 

\begin{enumerate}[leftmargin=*]
\item Given two rings $R$ and $S$ we let $\Hom_{\comRing}(R,S)$ be the set of ring homomorphisms from $R$ to $S$. When there is little room for confusion we will often drop the subscript $\comRing$ and simply write $\Hom(R,S)$. 
	\begin{enumerate}[label=(\alph*),ref=\theenumi\alph*]
		\item Prove or disprove: If $R$ and $S$ are rings, then $\Hom(R,S)$ is a ring under point-wise multiplication and point-wise addition. 
		\item\label{ex-part:hom-set-exm} Give an explicit description of the hom-sets below:
		\begin{multicols}{2}
		\begin{enumerate}[(i)]
			\item $\Hom(\ZZ, \ZZ)$
			\item $\Hom(\ZZ/\langle 2 \rangle, \ZZ)$
			\item $\Hom(\ZZ, \ZZ/\langle 2 \rangle)$
			\item $\Hom(\ZZ, \CC[x])$
		\end{enumerate}
		\end{multicols}
		\item Based on part \ref{ex-part:hom-set-exm} state and prove a characterization of $\Hom(\ZZ,R)$ for any ring $R$. 
	\end{enumerate}
		
	\item Let $\phi:R \to S$ be a ring homomorphism.
	\begin{enumerate}[label=(\alph*),ref=\theenumi\alph*]
		\item Prove that if $J \subset S$ is an ideal then $\phi^{-1}(J) \subset R$ is an ideal.
		\item Assuming $\phi$ is injective, give a nice description of $\phi^{-1}(J)$ for an ideal $J \subset S$. 
		\item Find a counterexample showing that if $I \subset R$ is an ideal, then $\phi(I)$ need not be an ideal.
		\item Prove that if $\phi$ is surjective then $\phi(I)$ is an ideal of $S$. 
		\item\label{ex-pt:ideal-preimage-image} Let $I \subset R$ be an ideal. Prove that if $\ker(\phi) \subset I$, then $\phi^{-1}(\phi(I))=I$.
		\item Is the assumption that $\ker(\phi) \subset I$ needed in part \ref{ex-pt:ideal-preimage-image}?
	\end{enumerate}
	
\end{enumerate}

\hrulefill 

As a small historical remark, despite it being common practice to include the existence of a multiplicative identity in the definition of a ring, in her original definition of a ring Emmy Noether did not do this. In E. Noether's 1921, article \textit{Idealtheorie in Ringbereichen} -- which can reasonably be viewed as the birth of commutative algebra -- she says, "\textit{Aus diesen Eigenschaften folgt die Existenz der Null; ein Ring braucht aber keine Einheit zu besitzen...}"; roughly translating to, "From these properties the existence of the zero element follows; however a ring is not required to possess a unit". 

\begin{center}
\includegraphics[scale=.3]{noether}
\end{center}

\hrulefill 
\end{document}





